\begin{abstract}
  \vspace*{1em}
  Face à la prolifération des malwares polymorphes, les signatures cryptographiques
classiques telles que MD5 ou SHA-256 sont structurellement obsolètes : un seul bit
modifié suffit à invalider toute comparaison exacte.
Ce rapport technique présente
\textbf{Binary-JEPA}, une architecture de hachage sémantique de binaires ELF64 fondée
sur l'apprentissage auto-supervisé, conçue pour produire une empreinte stable face aux
transformations compilatoires et, à terme, aux techniques d'obfuscation modernes.

Notre approche repose sur trois contributions articulées. Premièrement, un pipeline
d'extraction statique traduisant le code binaire en représentation intermédiaire VEX IR
via \texttt{angr}, puis modélisant la topologie du programme sous forme d'un
\textit{Bag-of-Paths} — un multiensemble de chemins d'exécution canonicalisés en
371~tokens abstraits par une projection sémantique stricte~$\Phi$, s'affranchissant des
artefacts compilatoires. Deuxièmement, une adaptation unidimensionnelle du cadre
I-JEPA~\cite{assran2023self} à l'analyse de séquences symboliques : un encodeur
convolutif léger est entraîné à prédire la représentation latente de positions masquées
plutôt qu'à reconstruire le token exact, forçant l'émergence d'une sémantique
comportementale sans supervision. Troisièmement, un paquet Python distribuable
(\texttt{elf64ctx}) encapsulant le pipeline complet avec téléchargement automatique
des artefacts depuis HuggingFace Hub.

Une preuve de concept qualitative dénote deux propriétés
attendues de l'espace latent : une invariance partielle à la compilation pour les paires
positives, et une séparabilité inter-programmes pour les
paires négatives, avec un effet super-additif lorsque les deux sources de divergence se cumulent. L'analyse de trois runs
d'entraînement révèle que la stratégie de masquage constitue le facteur de convergence
dominant : seul le masquage aléatoire produit un plateau stable à
$\mathcal{L} \approx 0.030$ sur 10~époques.

Ce rapport constitue un \textit{working paper} : les hypothèses de résilience face aux
obfuscations restent à valider empiriquement, et l'évaluation quantitative fait l'objet des travaux en cours. 
\end{abstract}

\keywords{Malware Analysis \and Semantic Hashing \and Self-Supervised Learning \and VEX IR \and \gls{bagofpaths} \and \gls{jepa}}
